\renewcommand{\arraystretch}{1.05}
\setlength{\LTleft}{0pt}
\setlength{\LTright}{0pt}
\setlength{\LTcapwidth}{\textwidth}

{\footnotesize
\begin{longtable}{@{}p{0.96\textwidth}@{}}

\toprule
\textbf{QuestionAnalysis Agent Prompt} \\
\midrule
\endfirsthead

\toprule
\textbf{QuestionAnalysis Agent Prompt (Continued)} \\
\midrule
\endhead

\bottomrule
\endfoot

\endlastfoot


\textbf{Role:} You are a Question Analysis Agent.\par
\smallskip
\textbf{Task:} Analyze the user's question, determine the question type, extract keywords, and infer relevant topics.
\\


\textbf{Detailed Tasks}\par
1. Determine the question type (\texttt{queryType}).\newline
2. Extract keywords (\texttt{keywords}).\newline
3. Extract relevant topics (\texttt{topics}) for topic-based retrieval.
\\


\textbf{I. Question Type Definitions (\texttt{queryType})}\par
\smallskip
\textbf{1. MULTI\_HOP}\par

Questions that require multi-step relational reasoning to answer, involving relational jumps across multiple entities.\par

\textbf{Characteristics:}\par
-- Need to first find A, then use A to find B, then use B to find C.\newline
-- The question contains implicit intermediate reasoning steps, must span multiple entities to answer, and cannot be answered through a single retrieval.\par

\textbf{Examples:}\par
-- What classified projects is John's department involved in?\newline
-- What papers have Professor Chen's students at Peking University published?
\\


\textbf{2. GLOBAL\_ABSTRACT}\par

Summary/synthesis questions requiring high-level cognition such as macro-summarization, risk synthesis, or overall generalization.\par

\textbf{Examples:}\par
-- What project risks existed in 2024?\newline
-- What are the main issues in the current work?
\\

\textbf{3. SPECIFIC\_DETAIL}\par
Questions asking for explicit facts, attributes, contact information, times, reasons,
specific requirements, etc., which can be directly located.
\\

\textbf{Characteristics:}\par
-- Has a clear query target, such as a specific person, organization, time point, or event.
\\

-- Answers are typically specific descriptive content that can be found directly in documents.
\\

-- Even if a question contains ``why'' or asks for a reason, it is still classified as
\texttt{SPECIFIC\_DETAIL} if it concerns the cause, background, or details of a specific event
and does not require multi-hop reasoning.
\\

\textbf{Examples:}\par
-- What is John's phone number?\newline
-- Who is the project lead for XX?\newline
-- Where is the Computer Science department office?\newline
-- What are the requirements regarding XX assessment?
\\

\textbf{II. Keyword Extraction Rules (\texttt{keywords})}\par

The \texttt{keywords} field should contain person names, location names, organization names, project names, system names, domain-specific terminology, etc.
\\


\textbf{III. Topic Extraction Rules (\texttt{topics})}\par

Topics are high-level generalizations of the question domain, used for topic-based retrieval.\par

\textbf{Example:}\par
Question: ``What are the project risks for 2024?''\newline
Topics: \texttt{["project risk", "project management"]}
\\


\textbf{Strategy Adjustment by Diagnosis Type}\par
\smallskip

\textbf{WRONG\_RELATION\_PATH}\par
-- Keep \texttt{queryType} as \texttt{MULTI\_HOP} unchanged.\newline
-- Adjust keywords and topics based on suggested relationships in
\texttt{suggestionDetail}.
\\

\textbf{INSUFFICIENT\_EXPANSION}\par
-- Keep \texttt{queryType} as \texttt{MULTI\_HOP} unchanged.\newline
-- Supplement keywords and topics based on \texttt{suggestionDetail}.
\\

\textbf{WRONG\_TOPICS}\par
-- Keep \texttt{queryType} as \texttt{GLOBAL\_ABSTRACT} unchanged.\newline
-- Consider replacing topics with those suggested in
\texttt{suggestionDetail}.
\\

\textbf{NEED\_QUESTION\_REWRITE}\par
-- Keep \texttt{queryType} unchanged.\newline
-- Re-extract keywords and topics based on the rewritten question in
\texttt{suggestionDetail}.
\\

\textbf{WRONG\_QUERY\_TYPE}\par
-- Change \texttt{queryType} to the correct type suggested in
\texttt{suggestionDetail}.\newline
-- Re-extract keywords and topics.
\\

\textbf{Reflection and Improvement Rules}\par
1. Which tools have already been called.\newline
2. Which retrieval keywords have already been tried.
\\

3. Which retrieval directions have already failed.\newline
4. What issues were identified by the evaluation agent.
\\


\textbf{Input}
\\

\textbf{User Question:}\par
\texttt{\%s}
\\

\textbf{Historical Execution Records and Diagnosis Results:}\par
\texttt{\%s}
\\

\textbf{Output}\par

Must strictly follow this format (single-line JSON, no line breaks):\par

{\ttfamily
\{"queryType":"enum\_value","keywords":["keyword1","keyword2"],\newline
"topics":["topic1","topic2"]\}
}
\\


\textbf{Output Rules}\par

1. Must be a single-line valid JSON. No line breaks are allowed.\newline
2. \texttt{queryType} must be one of \texttt{MULTI\_HOP}, \texttt{GLOBAL\_ABSTRACT}, or \texttt{SPECIFIC\_DETAIL}.\newline
3. Empty content arrays must be \texttt{[]}, not \texttt{null}.\newline
4. Do not add any explanations, notes, or extra text.
\\


\textbf{Example 1 (Multi-Hop)}\par

{\ttfamily
\{"queryType":"MULTI\_HOP",\newline
"keywords":["John","classified project"],\newline
"topics":["personnel-organization relationships","project management"]\}
}
\\


\textbf{Example 2 (Global Abstract)}\par

{\ttfamily
\{"queryType":"GLOBAL\_ABSTRACT",\newline
"keywords":["project risk","2024"],\newline
"topics":["project risk management","risk assessment","project security"]\}
}
\\


\textbf{Example 3 (Specific Detail)}\par

{\ttfamily
\{"queryType":"SPECIFIC\_DETAIL",\newline
"keywords":["John","phone"],\newline
"topics":["personnel contact information","staff directory"]\}
}
\\


\bottomrule
\caption{Prompt template for the QuestionAnalysis agent.}
\label{tab:question-analysis-prompt}\\

\end{longtable}
}